\documentclass[UTF8,a4paper,11pt]{ctexart}

\usepackage{amsmath,amssymb,bm,mathtools}
\usepackage{booktabs}
\usepackage{enumitem}
\usepackage{geometry}
\usepackage{xurl}
\usepackage[colorlinks=true,linkcolor=blue,citecolor=blue,urlcolor=blue]{hyperref}

\geometry{left=2.5cm,right=2.5cm,top=2.6cm,bottom=2.6cm}
\setlist{nosep}
\linespread{1.12}

\newcommand{\dd}{\mathrm{d}}
\newcommand{\jj}{\mathrm{j}}
\newcommand{\KK}{\mathcal{K}}
\newcommand{\RR}{\mathbb{R}}
\newcommand{\CC}{\mathbb{C}}
\newcommand{\cA}{\mathcal{A}}
\newcommand{\cB}{\mathcal{B}}
\newcommand{\cH}{\mathcal{H}}
\newcommand{\cI}{\mathcal{I}}
\newcommand{\cK}{\mathcal{K}}
\newcommand{\cL}{\mathcal{L}}
\newcommand{\cM}{\mathcal{M}}
\newcommand{\cO}{\mathcal{O}}
\newcommand{\transpose}{\mathsf{T}}
\newcommand{\herm}{\mathsf{H}}

\title{ALPS 快速扫频方法的数学模型推导}
\author{独立推导说明}
\date{\today}

\begin{document}

\maketitle

\begin{abstract}
ALPS（Adaptive Lanczos--Padé Sweep，自适应 Lanczos--Padé 扫频）是计算电磁学中常用的快速频率扫描思想：在少量中心频率处进行完整电磁求解，然后用由 Lanczos Krylov 子空间生成的 Padé 有理逼近在一段频带内外推场量和端口响应。本文不调用、不摘录、不依赖本工程已有文档，而是从频域 Maxwell 方程、有限维离散系统、矩匹配、Lanczos 投影和自适应频带覆盖几个层次独立推导其数学模型。推导中的 ALPS 指公开数值线性代数意义下的 Adaptive Lanczos--Padé Sweep；不同商业软件的工程实现可能包含额外的网格、端口和误差控制细节。
\end{abstract}

\tableofcontents

\section{扫频问题的抽象形式}

设时间因子为 $\exp(st)$，实际频率轴上 $s=\jj\omega$。计算电磁学中的频率扫描要求在频带
\[
    \Omega=[\omega_{\min},\omega_{\max}]
\]
上求解场量、端口阻抗、导纳或散射参数。若对每个频点都重新装配并分解大型线性系统，成本通常随频点数线性增长。ALPS 的目标是把大量频点的完整求解替换为少量中心频率处的完整求解，加上一组低阶有理模型的快速求值。

离散后的频域电磁问题可统一写成
\begin{equation}
    A(s)x(s)=b(s),\qquad y(s)=C(s)^{\herm}x(s)+d(s),
    \label{eq:full-system}
\end{equation}
其中 $x(s)\in\CC^N$ 是有限元、边界元或矩量法未知量，$y(s)\in\CC^r$ 是端口响应或观测量。快速扫频真正要近似的是输入到输出的频率响应
\begin{equation}
    H(s)=C(s)^{\herm}A(s)^{-1}B(s)+D(s),
    \label{eq:transfer-general}
\end{equation}
而不一定需要在每个频点重建全空间场解。

\section{由 Maxwell 方程得到频率参数线性系统}

在各向同性线性介质中，频域 Maxwell 方程可写为
\begin{align}
    \nabla\times E &= -s\mu H-M_s,\\
    \nabla\times H &= (\sigma_e+s\varepsilon)E+J_s,
\end{align}
其中 $\varepsilon,\mu,\sigma_e$ 分别为介电常数、磁导率和电导率。为突出 ALPS 的数学结构，先忽略磁流源 $M_s$。消去磁场 $H$ 得到电场的 curl--curl 方程
\begin{equation}
    \nabla\times\mu^{-1}\nabla\times E
    +s\sigma_e E+s^2\varepsilon E
    =-sJ_s .
    \label{eq:curlcurl}
\end{equation}
在实际频率轴 $s=\jj\omega$ 上，式 \eqref{eq:curlcurl} 变为常见形式
\[
    \nabla\times\mu^{-1}\nabla\times E
    +\jj\omega\sigma_e E-\omega^2\varepsilon E
    =-\jj\omega J_s .
\]

令 $V_h=\operatorname{span}\{N_1,\ldots,N_N\}$ 为边元有限元空间，取
\[
    E_h(s)=\sum_{n=1}^N x_n(s)N_n .
\]
对测试函数 $F_h\in V_h$ 作 Galerkin 离散，可得二次参数矩阵系统
\begin{equation}
    \bigl(K+sD+s^2M+A_{\Gamma}(s)\bigr)x(s)=b(s),
    \label{eq:fem-quadratic}
\end{equation}
其中
\begin{align}
    K_{mn} &=
    \int_{\Omega_e}(\nabla\times N_n)\cdot\mu^{-1}(\nabla\times\overline{N_m})\,\dd V,\\
    D_{mn} &=
    \int_{\Omega_e}N_n\cdot\sigma_e\overline{N_m}\,\dd V,\\
    M_{mn} &=
    \int_{\Omega_e}N_n\cdot\varepsilon\overline{N_m}\,\dd V .
\end{align}
$A_{\Gamma}(s)$ 表示辐射边界、吸收边界、PML 或波端口所产生的频率相关边界项。若边界项在局部频带内解析，它可以与 $K+sD+s^2M$ 一起进入后续 Taylor 展开或有理近似。

矩量法或混合位势积分方程也可得到同类模型。例如对电流基函数 $f_n$ 展开表面电流 $J_h=\sum_n i_n f_n$，常见阻抗矩阵可抽象写为
\begin{equation}
    Z(s)i(s)=v(s),
    \label{eq:mom-system}
\end{equation}
其中 $Z_{mn}(s)$ 来自含有 Helmholtz Green 函数
\[
    G_k(r,r')=\frac{\exp(-\jj k|r-r'|)}{4\pi |r-r'|},
    \qquad k=\omega\sqrt{\mu\varepsilon},
\]
的双重积分。只要选定中心频率附近没有材料分支点或系统奇点，$Z(s)$ 同样可看作 $s$ 的解析矩阵函数。因此，无论来自有限元、边界元还是矩量法，ALPS 都可从统一的参数系统 \eqref{eq:full-system} 出发。

\section{中心频率展开与一阶线性化}

取中心频率 $s_0=\jj\omega_0$，令
\[
    \sigma=s-s_0 .
\]
若系统已经是仿射形式
\begin{equation}
    A(s)=A_0+\sigma A_1,\qquad b(s)=b_0,
    \label{eq:affine-system}
\end{equation}
则 ALPS 的核心推导最为直接。对于二次有限元系统 \eqref{eq:fem-quadratic}，也可以精确地转化为一阶仿射 pencil。将
\[
    K+sD+s^2M
    =\widetilde K+\sigma\widetilde D+\sigma^2M
\]
其中
\[
    \widetilde K=K+s_0D+s_0^2M,\qquad
    \widetilde D=D+2s_0M .
\]
引入增广未知量
\[
    z(\sigma)=
    \begin{bmatrix}
        x(\sigma)\\
        \sigma x(\sigma)
    \end{bmatrix},
\]
则二次系统等价于
\begin{equation}
    \left(
    \underbrace{
    \begin{bmatrix}
        0&-I\\
        \widetilde K&\widetilde D
    \end{bmatrix}}_{\cA_0}
    +\sigma
    \underbrace{
    \begin{bmatrix}
        I&0\\
        0&M
    \end{bmatrix}}_{\cA_1}
    \right)
    \begin{bmatrix}
        x\\
        \sigma x
    \end{bmatrix}
    =
    \begin{bmatrix}
        0\\
        b
    \end{bmatrix}.
    \label{eq:quadratic-linearization}
\end{equation}
第一行给出 $\sigma x-\sigma x=0$，第二行正是
$\widetilde Kx+\sigma\widetilde D x+\sigma^2Mx=b$。更高阶多项式矩阵
\[
    A(s_0+\sigma)=A_0+\sigma A_1+\cdots+\sigma^pA_p
\]
也可通过伴随线性化转化为一阶仿射形式，未知量为
$[x,\sigma x,\ldots,\sigma^{p-1}x]^\transpose$。一般解析矩阵 $A(s)$ 可先在 $s_0$ 附近截断 Taylor 展开，再通过原始残差验证近似有效性。

因此，后续只需研究标准形式
\begin{equation}
    (\cA_0+\sigma\cA_1)u(\sigma)=\cB,
    \qquad
    y(\sigma)=\cL^{\herm}u(\sigma),
    \label{eq:standard-pencil}
\end{equation}
其中 $\cA_0$ 是中心频率处需要分解的大矩阵。若 $\cA_0$ 非奇异，则
\begin{equation}
    u(\sigma)
    =
    (\cA_0+\sigma\cA_1)^{-1}\cB
    =
    (I-\sigma\cH)^{-1}R,
    \label{eq:resolvent}
\end{equation}
其中
\begin{equation}
    R=\cA_0^{-1}\cB,\qquad
    \cH=-\cA_0^{-1}\cA_1 .
    \label{eq:H-R-definition}
\end{equation}
式 \eqref{eq:resolvent} 说明：一次完整频点求解或分解 $\cA_0$ 后，局部频率响应变成了矩阵 resolvent 的有理逼近问题。

\section{矩展开与 Padé 逼近}

当 $|\sigma|$ 小于 resolvent 收敛半径时，
\begin{equation}
    (I-\sigma\cH)^{-1}
    =
    \sum_{n=0}^{\infty}\sigma^n\cH^n .
\end{equation}
于是单输入单输出响应
\begin{equation}
    h(\sigma)=\ell^{\herm}u(\sigma)
\end{equation}
具有幂级数展开
\begin{equation}
    h(\sigma)
    =
    \sum_{n=0}^{\infty}m_n\sigma^n,
    \qquad
    m_n=\ell^{\herm}\cH^n r,
    \label{eq:moments}
\end{equation}
其中 $r=\cA_0^{-1}b$。$m_n$ 称为频率响应在中心点 $s_0$ 的矩。

Padé 逼近用有理函数代替多项式 Taylor 截断。给定整数 $p,q$，寻找
\begin{equation}
    h_{[p/q]}(\sigma)
    =
    \frac{P_p(\sigma)}{Q_q(\sigma)}
    =
    \frac{a_0+a_1\sigma+\cdots+a_p\sigma^p}
    {1+c_1\sigma+\cdots+c_q\sigma^q}
    \label{eq:pade}
\end{equation}
使
\begin{equation}
    Q_q(\sigma)h(\sigma)-P_p(\sigma)
    =\cO(\sigma^{p+q+1}).
    \label{eq:pade-condition}
\end{equation}
将 $h(\sigma)=\sum_{n\ge0}m_n\sigma^n$ 代入，可得分母系数满足
\begin{equation}
    \sum_{i=0}^{q}c_i m_{n-i}=0,
    \qquad
    n=p+1,\ldots,p+q,
    \qquad c_0=1,
    \label{eq:pade-den}
\end{equation}
分子系数为
\begin{equation}
    a_n=\sum_{i=0}^{\min(n,q)}c_i m_{n-i},
    \qquad n=0,\ldots,p .
    \label{eq:pade-num}
\end{equation}
与 Taylor 多项式相比，Padé 逼近允许分母出现近似极点，因此在谐振、传输零点或缓慢衰减模态附近通常能用更低阶数描述更宽频带。

直接使用 \eqref{eq:moments}--\eqref{eq:pade-num} 会遇到两个困难：一是高阶矩需要反复求解 $\cA_0$，二是矩序列容易严重病态。ALPS 的 ``Lanczos--Padé'' 部分正是用 Lanczos 投影隐式生成这些矩匹配条件。

\section{Lanczos 投影产生 Padé 有理模型}

先考虑单输入单输出情形。构造右 Krylov 子空间和左 Krylov 子空间
\begin{align}
    \cK_m(\cH,r)
    &=\operatorname{span}\{r,\cH r,\ldots,\cH^{m-1}r\},\\
    \cK_m(\cH^{\herm},\ell)
    &=\operatorname{span}\{\ell,\cH^{\herm}\ell,\ldots,(\cH^{\herm})^{m-1}\ell\}.
\end{align}
双边 Lanczos 过程生成矩阵
\[
    V_m=[v_1,\ldots,v_m],\qquad
    W_m=[w_1,\ldots,w_m],
\]
满足
\begin{equation}
    W_m^{\herm}V_m=I_m,\qquad
    T_m=W_m^{\herm}\cH V_m ,
    \label{eq:lanczos-proj}
\end{equation}
其中 $T_m$ 在无 breakdown 时为三对角矩阵。用 Petrov--Galerkin 条件近似
\[
    u_m(\sigma)=V_m z_m(\sigma),
\]
并要求残差对左空间正交：
\[
    W_m^{\herm}\bigl((I-\sigma\cH)V_mz_m-R\bigr)=0 .
\]
由此得到低阶系统
\begin{equation}
    (I_m-\sigma T_m)z_m(\sigma)=W_m^{\herm}R .
    \label{eq:reduced-system}
\end{equation}
响应近似为
\begin{equation}
    h_m(\sigma)
    =
    \ell^{\herm}V_m
    (I_m-\sigma T_m)^{-1}
    W_m^{\herm}r .
    \label{eq:lanczos-pade-response}
\end{equation}
若起始向量按 $r=\beta v_1$、$\ell=\eta w_1$ 归一化，则
\begin{equation}
    h_m(\sigma)
    =
    \eta^*\beta\,
    e_1^{\transpose}
    (I_m-\sigma T_m)^{-1}
    e_1 .
    \label{eq:scalar-lanczos-pade}
\end{equation}
式 \eqref{eq:scalar-lanczos-pade} 已经是一个分母为
\[
    Q_m(\sigma)=\det(I_m-\sigma T_m)
\]
的有理函数；其极点由 $T_m$ 的 Ritz 值决定，即
\[
    \sigma_k\approx\lambda_k(T_m)^{-1},
    \qquad
    s_k\approx s_0+\lambda_k(T_m)^{-1}.
\]
这些极点近似原始电磁系统在中心点附近的谐振或自然模态。

Lanczos--Padé 的关键矩匹配性质可以写成
\begin{equation}
    \ell^{\herm}\cH^n r
    =
    \ell^{\herm}V_mT_m^nW_m^{\herm}r,
    \qquad
    n=0,1,\ldots,2m-1,
    \label{eq:moment-matching}
\end{equation}
在无 breakdown 且双正交条件保持良好时成立。将 \eqref{eq:lanczos-pade-response} 展开为
\[
    h_m(\sigma)
    =
    \sum_{n=0}^{\infty}
    \ell^{\herm}V_mT_m^nW_m^{\herm}r\,\sigma^n,
\]
再与 \eqref{eq:moments} 比较，可知
\begin{equation}
    h(\sigma)-h_m(\sigma)=\cO(\sigma^{2m}).
\end{equation}
因此 $m$ 步双边 Lanczos 给出的低阶模型等价于一个典型的 $[m-1/m]$ Padé 逼近，但避免显式计算高阶矩和求解病态 Hankel 方程。

\section{多端口与矩阵有理模型}

电磁器件常有多个端口或多个激励。设
\[
    (\cA_0+\sigma\cA_1)U(\sigma)=\cB,\qquad
    Y(\sigma)=\cL^{\herm}U(\sigma),
\]
其中 $\cB\in\CC^{N\times n_i}$，$\cL\in\CC^{N\times n_o}$。由 \eqref{eq:H-R-definition} 得
\begin{equation}
    Y(\sigma)
    =
    \cL^{\herm}(I-\sigma\cH)^{-1}R,
    \qquad R=\cA_0^{-1}\cB .
\end{equation}
块 Lanczos 构造
\[
    \operatorname{range}(V_m)
    =
    \cK_m(\cH,R),
    \qquad
    \operatorname{range}(W_m)
    =
    \cK_m(\cH^{\herm},\cL),
\]
并保持 $W_m^{\herm}V_m=I$。于是矩阵响应的低阶模型为
\begin{equation}
    Y_m(\sigma)
    =
    \cL^{\herm}V_m
    (I-\sigma T_m)^{-1}
    W_m^{\herm}R,
    \qquad
    T_m=W_m^{\herm}\cH V_m .
    \label{eq:block-rom}
\end{equation}
该模型同时匹配多个输入输出通道的块矩
\[
    M_n=\cL^{\herm}\cH^nR .
\]
实际软件中常以端口电压、电流或功率波作为输入输出，先得到阻抗矩阵 $Z_m(s)$ 或导纳矩阵 $Y_m(s)$，再转换为散射矩阵。例如对参考阻抗矩阵 $Z_0$，
\begin{equation}
    S_m(s)
    =
    \bigl(Z_m(s)-Z_0\bigr)
    \bigl(Z_m(s)+Z_0\bigr)^{-1},
    \label{eq:s-param}
\end{equation}
或使用相应的功率归一化形式。由于 $Z_m(s)$ 是低阶有理矩阵，$S_m(s)$ 也是可快速求值的矩阵有理函数。

\section{自适应扫频策略}

单个中心频率的 Padé 模型只在某个邻域内可靠。ALPS 的 ``Adaptive'' 部分就是自动选择中心频率、模型阶数和有效频带。给定误差容限 $\tau$，可定义原始方程残差指标
\begin{equation}
    \varepsilon_A(s)
    =
    \frac{\|b(s)-A(s)x_m(s)\|_2}
         {\|b(s)\|_2+\|A(s)\|_2\|x_m(s)\|_2},
    \label{eq:residual-error}
\end{equation}
以及响应收敛指标
\begin{equation}
    \varepsilon_Y(s)
    =
    \frac{\|Y_m(s)-Y_{m-\Delta m}(s)\|_F}
         {\|Y_m(s)\|_F+\epsilon_{\mathrm{mach}}}.
    \label{eq:response-error}
\end{equation}
若 $A(s)$ 非奇异，则残差给出误差关系
\begin{equation}
    x(s)-x_m(s)=A(s)^{-1}\bigl(b(s)-A(s)x_m(s)\bigr),
\end{equation}
从而
\begin{equation}
    \frac{\|x(s)-x_m(s)\|_2}{\|x(s)\|_2}
    \le
    \kappa_2(A(s))
    \frac{\|b(s)-A(s)x_m(s)\|_2}{\|b(s)\|_2},
    \label{eq:error-bound}
\end{equation}
其中 $\kappa_2(A)$ 为二范数条件数。靠近谐振或传输零点时 $\kappa_2(A)$ 可能很大，因此通常同时使用残差、响应变化和采样点校验。

围绕中心 $\omega_c$ 的有效半径可定义为
\begin{equation}
    \rho_c
    =
    \sup\left\{\rho>0:
    \max_{\omega\in[\omega_c-\rho,\omega_c+\rho]\cap\Omega}
    \max(\varepsilon_A(\jj\omega),\varepsilon_Y(\jj\omega))
    \le\tau
    \right\}.
    \label{eq:valid-radius}
\end{equation}
实际算法可描述如下：
\begin{enumerate}
    \item 选取尚未覆盖频段中的中心频率 $\omega_c$，装配并分解 $\cA_0=\cA(s_c)$。
    \item 用 $\cA_0^{-1}$ 作为线性求解器，执行 $m$ 步 Lanczos 或块 Lanczos，得到 $T_m,V_m,W_m$。
    \item 在候选采样点上计算低阶模型 \eqref{eq:block-rom}，并评估 \eqref{eq:residual-error}--\eqref{eq:response-error}。
    \item 若误差满足容限，则接受由 \eqref{eq:valid-radius} 给出的局部频带；若不满足，则增加 $m$ 或缩小局部频带。
    \item 对未覆盖的频段重复上述过程，直到 $\Omega$ 被若干重叠局部模型覆盖。
\end{enumerate}
最终扫频结果可表示为
\begin{equation}
    Y_{\mathrm{ALPS}}(s)=Y_{m_k}^{(k)}(s),
    \qquad
    s\in \cI_k,
\end{equation}
其中 $\{\cI_k\}$ 为各中心频率模型的有效子区间。若多个区间重叠，常选择离当前频率最近的中心模型，或用残差较小的模型。

\section{与 AWE 和直接扫频的关系}

渐近波形估计（AWE）同样从中心频率的矩展开出发，但通常显式计算矩并求解 Padé 系数方程。ALPS 可看作 AWE 的 Krylov/Lanczos 稳定化版本：
\begin{itemize}
    \item 直接扫频：每个频点求解一次全阶系统，鲁棒但昂贵。
    \item AWE：在中心点计算矩并构造 Padé 模型，频点求值便宜，但高阶矩容易病态。
    \item ALPS：用 Lanczos 子空间隐式实现 Padé 矩匹配，并用自适应中心覆盖宽频段。
\end{itemize}
因此，ALPS 的数学本质不是改变 Maxwell 方程，而是对离散后的频率参数线性系统进行局部有理降阶。

\section{适用条件与数值注意事项}

\begin{enumerate}
    \item \textbf{解析性。} 在每个中心频率邻域内，$A(s)$、$b(s)$ 和输出算子应为解析或可由低阶解析模型近似。材料强色散、模式截止点和非光滑端口模型会缩小有效半径。
    \item \textbf{非奇异性。} 中心矩阵 $\cA_0$ 必须可分解。若中心点过于接近自然谐振，Lanczos 模型可能需要更高阶数，残差也会放大。
    \item \textbf{双正交稳定性。} 双边 Lanczos 可能出现 breakdown 或双正交性丢失。工程实现通常需要再正交、块 Lanczos、look-ahead Lanczos 或 Arnoldi 变体。
    \item \textbf{被动性。} 对无源电磁网络，物理极点应满足稳定性要求。数值降阶可能产生轻微非被动性，必要时需做被动性检查或修正。
    \item \textbf{误差验证。} Padé 模型阶数较低时外推很快，但不能无限延伸。自适应分段和原始方程残差是宽频可靠性的关键。
\end{enumerate}

\section{结论}

ALPS 快速扫频的核心数学模型可以概括为
\[
    A(s)x(s)=b(s)
    \quad\Longrightarrow\quad
    u(\sigma)=(I-\sigma\cH)^{-1}R
    \quad\Longrightarrow\quad
    y_m(\sigma)=
    \cL^{\herm}V_m(I-\sigma T_m)^{-1}W_m^{\herm}R .
\]
其中第一步来自 Maxwell 方程离散化，第二步来自中心频率处的一阶线性化和矩阵分解，第三步来自 Lanczos Krylov 投影。由于 $y_m$ 是 Padé 型有理函数，它能以少量自由度表示谐振电磁响应；由于有效频带由残差和响应收敛自适应决定，它又能覆盖宽频扫频任务。这就是 ALPS 相比逐点求解能够显著减少完整电磁求解次数的根本原因。

\begin{thebibliography}{9}

\bibitem{feldmann1995}
P. Feldmann and R. W. Freund,
``Efficient linear circuit analysis by Padé approximation via the Lanczos process,''
\emph{IEEE Transactions on Computer-Aided Design of Integrated Circuits and Systems},
vol. 14, no. 5, pp. 639--649, 1995.
\url{https://doi.org/10.1109/43.384428}

\bibitem{bai2001}
Z. Bai and R. W. Freund,
``A partial Padé-via-Lanczos method for reduced-order modeling,''
\emph{Linear Algebra and its Applications},
vol. 332--334, pp. 139--164, 2001.
\url{https://doi.org/10.1016/S0024-3795(00)00291-3}

\bibitem{bracken1999}
J. E. Bracken, D.-K. Sun, and Z. J. Cendes,
``Characterization of electromagnetic devices via reduced-order models,''
\emph{Computer Methods in Applied Mechanics and Engineering},
vol. 169, no. 3--4, pp. 311--330, 1999.
\url{https://doi.org/10.1016/S0045-7825(98)00160-1}

\bibitem{baker1996}
G. A. Baker and P. Graves-Morris,
\emph{Padé Approximants},
2nd ed., Cambridge University Press, 1996.

\bibitem{ansys-fast-sweep}
Ansys HFSS documentation,
``Fast Frequency Sweeps,''
accessed 2026.
\url{https://ansyshelp.ansys.com/public/Views/Secured/Electronics/v251/en/Subsystems/HFSS/Content/HFSS/FastFrequencySweeps.htm}

\end{thebibliography}

\end{document}
